% Trouver les mots manquants
\begin{itemize}
    \item Une force appliquée à un solide peut le \dotline{3cm}, le
          \dotline{3cm} ou augmenter sa \dotline{3cm}.
          
    \item Le travail d'une force dont la direction coïncide avec le déplacement
          de son point d'application s'obtient en multipliant la \dotline{3cm}
          de la force par la \dotline{3cm} du déplacement. Il est positif si la
          force et le déplacement ont le \dotline{3cm}~; il est \dotline{3cm}
          dans le cas contraire.
          
    \item Dans le Système International, les travaux de forces s'expriment en
          \dotline{3cm} (\dotline{1cm}) et la puissance d'une force en
          \dotline{3cm} (\dotline{1cm}).
\end{itemize}

